\par\medskip\noindent\SGAInlineRefTarget{sga3:xix:paragraph:1-8}\textbf{1.8.}\quad
A \(k\)-group \(G\) is said to be \emph{semisimple} if it is smooth,
affine, and connected, and if
\[
  \operatorname{rad}(G)=\{e\},
\]
that is, if \(G\) is reductive and
\(\underline{\operatorname{Centr}}(G)\) is finite.  If \(G\) is an
arbitrary smooth affine connected \(k\)-group, then
\(G/\operatorname{rad}(G)\) is semisimple (\emph{Bible}, §9.4,
Proposition 2), and \(G/\operatorname{rad}^{u}(G)\) is reductive.  The
\emph{semisimple rank of \(G\)}, denoted by
\(\operatorname{rgss}(G/k)\), or simply by \(\operatorname{rgss}(G)\), is
the reductive rank of \(G/\operatorname{rad}(G)\).
\SGARefTarget{sga3:xix:section:1:8}

Thus, if \(G\) is reductive,
\[
  \operatorname{rgred}(G)
  =
  \operatorname{rgss}(G)+\dim\bigl(\operatorname{rad}(G)\bigr).
\]

If \(G\) is a smooth affine connected \(k\)-group and \(Q\) is a central
subtorus of \(G\), then \(G/Q\) is semisimple if and only if \(G\) is
reductive and \(Q=\operatorname{rad}(G)\).  Indeed, if \(G/Q\) is
semisimple, then \(Q\supset\operatorname{rad}(G)\); hence
\(\operatorname{rad}(G)\) is a torus, so \(G\) is reductive, and
\(\operatorname{rad}(G)\) is the maximal torus of
\(\underline{\operatorname{Centr}}(G)\), whence
\(\operatorname{rad}(G)=Q\).  If \(G\) is reductive and \(Q\) is a central
subgroup, then \(G/Q\) is reductive and
\[
  \operatorname{rgss}(G)=\operatorname{rgss}(G/Q).
\]

\par\medskip\noindent\SGAInlineRefTarget{sga3:xix:paragraph:1-9}\textbf{1.9.}\quad
Let \(K\) be an algebraically closed extension of \(k\), and let \(G\) be
a smooth affine connected \(k\)-group.  Then \(G\) is reductive
(respectively, semisimple) if and only if \(G_K\) is, and
\[
  \operatorname{rgred}(G/k)=\operatorname{rgred}(G_K/K),
  \qquad
  \operatorname{rgss}(G/k)=\operatorname{rgss}(G_K/K).
\]
\SGARefTarget{sga3:xix:section:1:9}

\par\medskip\noindent\SGAInlineRefTarget{sga3:xix:paragraph:1-10}\textbf{1.10.}\quad
Let \(G\) be a smooth connected \(k\)-group and let \(T\) be a torus of
\(G\).\footnote{Editor's note: the following sentence has been added.}
Write \(\mathfrak g\) for the \(k\)-Lie algebra of \(G\), that is,
\(\mathfrak g=\operatorname{Lie}(G)=\operatorname{Lie}(G/k)(k)\), and
similarly write \(\mathfrak t=\operatorname{Lie}(T)\).  Under the action
of \(T\), through \(T\longrightarrow G\) and the adjoint representation
of \(G\), the Lie algebra \(\mathfrak g\) decomposes as
\begin{equation}
  \mathfrak g
  =
  \mathfrak g^{0}
  \oplus
  \bigoplus_{\alpha\in R}\mathfrak g^{\alpha},
  \tag{1.10.1}\label{sga3:xix:paragraph:1-10:equation:1-10-1}
\end{equation}
where the elements \(\alpha\in R\) are nontrivial characters of \(T\)
and the spaces \(\mathfrak g^{\alpha}\) are nonzero.  The characters
\(\alpha\in R\) are called the \emph{roots of \(G\) relative to \(T\)}.
By \SGARefLink{sga3:ii:theorem:5-2-3:item:i}{Exposé II, 5.2.3(i)},
\begin{equation}
  \mathfrak g^{0}
  =
  \operatorname{Lie}\bigl(\underline{\operatorname{Centr}}_{G}(T)\bigr).
  \tag{1.10.2}
\end{equation}
\SGARefTarget{sga3:xix:equation:1:10:1}
\SGARefTarget{sga3:xix:equation:1:10:2}

In particular,\footnote{Editor's note: the following argument has been
given in detail.} since
\(\underline{\operatorname{Centr}}_{G}(T)\) is connected (by \SGARefLink{sga3:xix:paragraph:1-3}{1.3} when
\(G\) is affine, and by \SGARefLink{sga3:xii:6.6:item:b}{Exposé XII, 6.6(b)}, in general), \(T\) is its own
centralizer if and only if
\(\mathfrak g^{0}=\operatorname{Lie}(T)\).  This condition implies that
\(T\) is maximal and that
\(\underline{\operatorname{Centr}}(G)\subset T\); hence, by \SGARefLink{sga3:xii:8.8:item:d:02}{Exposé XII,
8.8(d)},
\begin{equation}
  \underline{\operatorname{Centr}}(G)
  =
  \operatorname{Ker}\bigl(T\longrightarrow
    \operatorname{GL}(\mathfrak g)\bigr)
  =
  \bigcap_{\alpha\in R}\operatorname{Ker}(\alpha).
  \tag{1.10.3}
\end{equation}
\SGARefTarget{sga3:xix:equation:1:10:3}
Moreover, \(\underline{\operatorname{Centr}}(G)\) is then affine, and
therefore the morphism
\(G\longrightarrow G/\underline{\operatorname{Centr}}(G)\) is affine.
Since the latter group is affine (\SGARefLink{sga3:xii:6.1}{Exposé XII, 6.1}), \(G\) is affine.

When \(G\) is reductive and \(T\) is maximal, the roots in the preceding
sense coincide with the roots in the sense of \emph{Bible}, §12.2,
Definition 1.  The latter are indeed roots in the sense of this number
(\emph{Bible}, §13.2, Theorem 1(c)), and their number is
\(\dim(G)-\dim(T)\) (\emph{Bible}, §13.4, Corollary 2 to Theorem 3).
Moreover, if \(\alpha\) is a root, then so is \(-\alpha\)
(\emph{Bible}, §12.2, Corollary to Proposition 1).  As usual, we write
the group law of
\(\operatorname{Hom}_{k\text{-}\mathrm{gr}}
  (T,\mathbb G_{m,k})\)
indifferently additively or multiplicatively.  It follows that, for
reductive \(G\),
\begin{equation}
  \dim(G)-\operatorname{rgred}(G)=\operatorname{Card}(R)
  \tag{1.10.4}
\end{equation}
is always even.
\SGARefTarget{sga3:xix:equation:1:10:4}

The semisimple rank of the reductive group \(G\) is the rank of the
subset \(R\) of the \(\mathbb Q\)-vector space
\[
  \operatorname{Hom}_{k\text{-}\mathrm{gr}}
    (T,\mathbb G_{m,k})\otimes_{\mathbb Z}\mathbb Q
\]
(\emph{Bible}, §13.3, Theorem 2).
\SGARefTarget{sga3:xix:section:1:10}
